\BourbakiExercisesSection

\begin{enumerate}
\def\labelenumi{\arabic{enumi}。}
\item
  设 G 是一个群，H 是 G 的一个正规子群，使得 G/H 是有限 n 阶循环群。设 x 是 G 中的一个元素，其像 \(\bar{x}\) 在 G/H 中生成 G/H。设 \(\phi\) 是 H 的自同构 \(h \mapsto xhx^{-1}\)，并设 \(y = x^{n}\)；于是 \(y \in H\)。证明 \(\phi(y) = y\)，并且 \(\phi^{n}\) 是由 y 定义的 H 的内自同构。设 \(\tau\) 是 Z 在 H 上的作用，由 \((m, h) \mapsto \phi^{m}(h)\) 定义，并设 \(E = H \times_{\tau} Z\) 为相应的半直积。证明 E 的元素 \((y^{-1}, n)\) 生成 E 的一个中心子群 \(C_{y}\)，且商 \(E/C_{y}\) 与 G 同构。
\item
  证明 Z 的每个中心扩张都是平凡的。
\item
  定义扩张（参见§5，习题 10）
\end{enumerate}

\[
\mathbf {Z} / 2 \mathbf {Z} \times \mathbf {Z} / 2 \mathbf {Z} \rightarrow \mathfrak {A} _ {4} \rightarrow \mathfrak {A} _ {3}
\]

\[
\mathbf {Z} / 2 \mathbf {Z} \times \mathbf {Z} / 2 \mathbf {Z} \rightarrow \mathfrak {S} _ {4} \rightarrow \mathfrak {S} _ {3}.
\]

证明这些扩张都是非平凡的，第一个扩张有一个截面，而第二个没有。

¶ 4. (a) 设 G 是一个阶为 mn 的有限群，且存在 G 的一个正规子群 H，H 是 m 阶循环群，商群 G/H 是 n 阶循环群。证明 G 可由两个元素 a, b 生成，使得 \(a^{m} = e\)，\(b^{n} = a^{r}\)，\(bab^{-1} = a^{s}\)，其中 r 和 s 是两个整数，使得 \(r(s - 1)\) 和 \(s^{n} - 1\) 是 m 的倍数（取 a 为生成 H 的元素，b 为生成 G/H 的陪集中的一个元素；将元素 \(b^{h}a^{k}b^{-h}\) 表示为 a 的幂，并特别应用于 h = n, k = 1 和 h = 1, k = r 的情形）。

\begin{enumerate}
\def\labelenumi{(\alph{enumi})}
\setcounter{enumi}{1}
\tightlist
\item
  *反之，设 G(m, n, r, s) 是由如下表现定义的群
\end{enumerate}

\[
\langle a, b; a ^ {m} = e, b ^ {n} = a ^ {r}, b a b ^ {- 1} = a ^ {s} \rangle ,
\]

（其中 \(m\) 和 \(n\) 是两个整数 \(\geqslant 0\)，而 \(r\) 和 \(s\) 是任意整数（参见 § 7，第 6 号）；证明，如果 \(m, r(s - 1)\) 和 \(s^n - 1\) 并非全为零，则 \(G(m, n, r, s)\) 是一个有限群，其阶为 \(qn\)，其中 \(q\) 是 \(m, |r(s - 1)|\) 和 \(|s^n - 1|\) 的最大公约数；在这个群中，由 \(a\) 生成的子群 H 是一个 \(q\) 阶正规子群，且 G/H 是一个 \(n\) 阶循环群（证明 \(G(m, n, r, s)\) 的每个元素都可以写成 \(a^x b^y\) 的形式，其中 \(x\) 和 \(y\) 是两个整数，使得 \(0 \leqslant x \leqslant q - 1\)，\(0 \leqslant y \leqslant n - 1\)，且 \(G(m, n, r, s)\) 同构于由满足上述条件的整数有序对 \((x, y)\) 组成的群，其合成律为

\[
(x, y). (x ^ {\prime}, y ^ {\prime}) = \left\{ \begin{array}{l l} (x + x ^ {\prime} s ^ {y}, y + y ^ {\prime}) & \text { if } y + y ^ {\prime} \leqslant n - 1 \\ (x + x ^ {\prime} s ^ {y} + r, y + y ^ {\prime} - n) & \text { if } y + y ^ {\prime} \geqslant n \end{array} \right.
\]

右端的第一坐标是模 \(q\) 的和。考察 \(m = r(s - 1) = s^n - 1 = 0\) 的情形。

群 G(n, 2, 0, -1) 称为 2n 阶二面体群，记作 \(\mathbf{D}_{n}\)；群 G(4, 2, 2, -1) 是一个 8 阶群，称为四元数群，记作 \(\mathfrak{Q}\)。证明在 \(\mathfrak{Q}\) 中每个子群都是正规的，并且异于 \(\{e\}\) 的诸子群之交是一个异于 \(\{e\}\) 的子群。证明 \(\mathbf{D}_{4}\) 不同构于 \(\mathfrak{Q}.\) *

\begin{enumerate}
\def\labelenumi{\arabic{enumi}。}
\setcounter{enumi}{4}
\tightlist
\item
  设 F 是一个中心为 \(\{e\}\) 的群，A 是它的自同构群。借助同态 Int: \(F \rightarrow A\) 把 F 等同于 A 的一个子群；令 \(\Gamma = A/F\)。
\end{enumerate}

\begin{enumerate}
\def\labelenumi{(\alph{enumi})}
\item
  证明扩张 \(\mathbf{F} \to \mathbf{A} \to \Gamma\) 仅当 \(\Gamma = \{e\}\) 时才是平凡的（注意 \(\mathbf{F}\) 在 \(\mathbf{A}\) 中的中心化子是 \(\{e\}\)）。
\item
  假设 \(\Gamma = \{e\}\)。证明群 G 的每个以群 F 为核的扩张
\end{enumerate}

\[
\mathbf {F} \rightarrow \mathbf {E} \rightarrow \mathbf {G}
\]

都是平凡的。

\begin{enumerate}
\def\labelenumi{\arabic{enumi}。}
\setcounter{enumi}{5}
\tightlist
\item
  设 I 是一个集合；记 \(\mathbf{F} = \mathbf{Z}, \mathbf{E} = \mathbf{Z} \times (\mathbf{Z}/2\mathbf{Z})^{\mathrm{I}}\) 和 \(\mathbf{G} = \mathbf{Z}/2\mathbf{Z} \times (\mathbf{Z}/2\mathbf{Z})^{\mathrm{I}}\)。
\end{enumerate}

\begin{enumerate}
\def\labelenumi{(\alph{enumi})}
\item
  定义一个非平凡扩张 \(\mathbf{F} \xrightarrow{i} \mathbf{E} \xrightarrow{p} \mathbf{G}\)。
\item
  证明：若 I 是无限的，则 E 同构于 \(F \times G\)。
\end{enumerate}

\begin{enumerate}
\def\labelenumi{\arabic{enumi}。}
\setcounter{enumi}{6}
\tightlist
\item
  设 G 和 A 是两个群，其中 A 是交换的。设 \(\tau\) 是 G 到 A 的自同构群 \(\operatorname{Aut}(\mathbf{A})\) 的一个同态；若 \(g \in \mathbf{G}\)、\(a \in \mathbf{A}\)，我们记 \(^g a = \tau(g)(a)\)。G 到 A 的交叉同态是指满足如下条件的任意映射 \(\phi: \mathbf{G} \to \mathbf{A}\)：
\end{enumerate}

\[
\phi (g g ^ {\prime}) = \phi (g) + ^ {g} \phi (g ^ {\prime}),
\]

其中群 A 以加法记法书写。G 到 A 的交叉同态在加法下构成一个群，记为 Z(G, A)。

\begin{enumerate}
\def\labelenumi{(\alph{enumi})}
\item
  若 \(a \in \mathbf{A}\)，映射 \(g \mapsto {}^g a - a\) 记为 \(\theta_a\)。证明 \(a \mapsto \theta_a\) 是一个同态 \(\theta\)，它把 \(\mathbf{A}\) 映入 \(Z(\mathbf{G}, \mathbf{A})\)。\(\theta\) 的核是 \(\mathbf{A}\) 的由在 \(G\) 作用下不变的元素组成的子群。\(\theta\) 的像记为 \(B(\mathbf{G}, A)\)。
\item
  设 \(\mathbf{X} = \mathbf{A} \times_{\tau} \mathbf{G}\) 为 \(\mathbf{G}\) 与 \(\mathbf{A}\) 的半直积。若 \(\phi\) 是 \(\mathbf{G}\) 到 \(\mathbf{A}\) 的一个映射，证明 \(g \mapsto (\phi(g), g)\) 是 \(\mathbf{X}\) 的一个截面，当且仅当 \(\phi\) 属于 \(\mathbf{Z}(\mathbf{G}, \mathbf{A})\)。为使对应于 \(\phi_1, \phi_2 \in \mathbf{Z}(\mathbf{G}, \mathbf{A})\) 的截面相互之间被 \(\mathbf{A}\) 的一个元素共轭，必要且充分的条件是 \(\phi_1 \equiv \phi_2 \bmod. \mathbf{B}(\mathbf{G}, \mathbf{A})\)。
\item
  设 \(\mathbf{F} \xrightarrow{i} \mathbf{E} \xrightarrow{p} \mathbf{G}\) 是 \(\mathbf{G}\) 的一个扩张，其核为群 \(\mathbf{F}\)，后者的中心等于 \(\mathbf{A}\)。假设：若 \(y \in \mathbf{E}\) 且 \(x = p(y)\)，则
\end{enumerate}

\[
i (x f) = y. i (f). y ^ {- 1}
\]

对所有 \(f \in \mathbf{A}\) 成立。若 \(\phi \in \mathbf{Z}(\mathbf{G}, \mathbf{A})\)，设 \(u_{\phi}\) 为 \(\mathbf{E}\) 到 \(\mathbf{E}\) 的由公式 \(u_{\phi}(y) = i(\phi(p(y))) .y\) 给出的映射。证明 \(\phi \mapsto u_{\phi}\) 是 \(\mathbf{Z}(\mathbf{G}, \mathbf{A})\) 到 \(\operatorname{Aut}(\mathbf{E})\)（扩张 \(\mathbf{E}\) 的自同构群）上的一个同构。这个同构把 \(\mathbf{B}(\mathbf{G}, \mathbf{A})\) 映到由自同构 \(\operatorname{Int}_{\mathbf{E}}(x)\)（其中 \(x\) 遍历 \(\mathbf{A}\)）所成的群。

¶ 8. 使用上述习题的记号，C(G, A) 用于表示从 G 到 A 的映射群。G 通过以下公式作用于 C(G, A)

\[
\left(^ {g} \phi\right) \left(g ^ {\prime}\right) = ^ {g} \left(\phi \left(g ^ {- 1} g ^ {\prime}\right)\right).
\]

若 \(\sigma\) 表示 \(\mathbf{G}\) 到 \(\operatorname{Aut}(\mathbf{C}(\mathbf{G},\mathbf{A}))\) 的相应同态，则半直积 \(\mathbf{C}(\mathbf{G},\mathbf{A}) \times_{\sigma} \mathbf{G}\) 记为 \(\mathbf{E}_0\)。

\begin{enumerate}
\def\labelenumi{(\alph{enumi})}
\item
  设 \(\varepsilon: \mathbf{A} \to \mathbf{C}(\mathbf{G}, \mathbf{A})\) 为这样的映射：它把每个 \(a \in \mathbf{A}\) 对应到等于 \(a\) 的常值映射。证明 \(\varepsilon\) 是一个单射同态，且与 \(\mathbf{G}\) 的作用相容。
\item
  设 \(\mathbf{A} \xrightarrow{i} \mathbf{E} \xrightarrow{p} \mathbf{G}\) 是 \(\mathbf{G}\) 的一个以 \(\mathbf{A}\) 为核的扩张，使得 \(i(ga) = xi(a)x^{-1}\)（若 \(a \in \mathbf{A}, x \in \mathbf{E}\) 且 \(g = p(x)\)）。设 \(\rho\) 是一个映射 \(\mathbf{G} \to \mathbf{E}\)，使得 \(p \circ \rho = \operatorname{Id}_{\mathbf{G}}\)。对所有 \(x \in \mathbf{E}\)，设 \(\phi_x\) 是 \(\mathbf{G}\) 到 \(\mathbf{A}\) 的一个映射，使得
\end{enumerate}

\[
x \rho (p (x ^ {- 1}) g) = i (\phi_ {x} (g)) \rho (g) \quad \text { for   all } g \in G.
\]

证明

\[
\phi_ {x y} = \phi_ {x} + ^ {p (x)} \phi_ {y} \quad \text { if } x, y \in \mathbf {E}.
\]

推断映射 \(\Phi: \mathbf{E} \to \mathbf{E}_0\)（它由

\[
\Phi (x) = (\phi_ {x}, p (x))
\]

定义）是一个同态，且使以下图表交换：

\[
\begin{array}{c c c c c} \mathrm{A} & \stackrel {{i}} {{\longrightarrow}} & \mathrm{E} & \stackrel {{p}} {{\longrightarrow}} & \mathrm{G} \\ \Big \downarrow_ {\varepsilon} & & \Big \downarrow_ {\Phi} & & \Big \downarrow_ {\mathrm{Id} _ {G}} \\ \mathrm{C(G,A)} & \longrightarrow & \mathrm{E} _ {0} & \longrightarrow & \mathrm{G} \end{array} .
\]

\begin{enumerate}
\def\labelenumi{(\alph{enumi})}
\setcounter{enumi}{2}
\tightlist
\item
  A 上的 G-均值是满足以下两个条件的同态
\end{enumerate}

\[
m \colon \mathbf {C} (\mathbf {G}, \mathbf {A}) \rightarrow \mathbf {A}
\]

满足以下两个条件：

\[
(c _ {1}) m \circ \varepsilon = \mathrm{Id} _ {\mathrm{A}}
\]

\[
(c _ {2}) m (^ {g} \phi) = ^ {g} m (\phi) \quad \text { if } g \in \mathrm{G}, \phi \in \mathrm{C} (\mathrm{G}, \mathrm{A}).
\]

设 \(m\) 是 A 上的 G-均值。证明，如果 \(\phi\) 是交叉同态且 \(a = m(\phi)\)，则 \(\phi = \theta_{-a}\)（参见习题 7）。特别地，

\[
\mathbf {Z} (\mathbf {G}, \mathbf {A}) = \mathbf {B} (\mathbf {G}, \mathbf {A}).
\]

证明映射 \((\phi, g) \mapsto (m(\phi), g)\) 是 \(\mathbf{E}_0\) 到 \(\mathbf{A} \times_{\tau} \mathbf{G}\) 的一个同态。由此推出，每个扩张 \(\mathbf{E}\)（\(\mathbf{G}\) 的、以 \(\mathbf{A}\) 为核且满足 \((b)\) 中条件者）都同构于 \(\mathbf{A} \times_{\tau} \mathbf{G}\)（使用复合 \(\mathbf{E} \xrightarrow{\Phi} \mathbf{E}_0 \to \mathbf{A} \times_{\tau} \mathbf{G}\)），从而有一个截面。证明两个这样的截面可以通过 \(\mathbf{E}\) 的一个内自同构（由 \(\mathbf{A}\) 的一个元素所定义）相互变换（使用习题 7 以及以下事实：\(\mathbf{Z}(\mathbf{G}, \mathbf{A}) = \mathbf{B}(\mathbf{G}, \mathbf{A})\)）。

\begin{enumerate}
\def\labelenumi{(\alph{enumi})}
\setcounter{enumi}{3}
\tightlist
\item
  假设 G 是 n 阶有限群，且映射 \(\alpha \mapsto n\alpha\) 是 A 的一个自同构。对所有 \(\phi \in \mathrm{C}(\mathrm{G}, \mathrm{A})\)，设 \(m(\phi)\) 表示 A 中满足
\end{enumerate}

\[
n. m (\phi) = \sum_ {g \in G} \phi (g).
\]

的唯一元素。证明 m 是 A 上的一个 G-均值。

这特别适用于 A 是阶与 n 互素的有限群的情形。

¶ 9. 设 F → E → G 是有限群的一个扩张。假设没有素数同时整除 F 的阶和 G 的阶。

\begin{enumerate}
\def\labelenumi{(\alph{enumi})}
\item
  假设 F 是可解的。证明存在一个截面 \(s: G \rightarrow E\)，并且两个这样的截面可以通过 F 中的一个元素相互共轭。（对 F 的可解类进行归纳论证；当 F 交换时，使用前一习题。）
\item
  证明在不假设 F 可解的情况下，存在截面† s: G → E。（对 G 的阶进行归纳论证。如果 p 整除 F 的阶，选择 F 的一个西罗 p-子群 P，并考虑它在 E 中的正规化子 N。N 在 G 中的像等于 G，参见习题 25。如果 N ≠ E，利用归纳假设得出结论；如果 N = E，使用 (a) 以及对 F/P → E/P → G 应用的归纳假设。）
\end{enumerate}

¶ 10. 设 G 是阶为 mn 的可解有限群，其中 m 和 n 没有公共素因子。证明存在 G 的一个阶为 m 的子群 H，并且 G 的每个阶整除 m 的子群都包含在 H 的某个共轭中（“霍尔定理”）。

（对 G 的阶进行归纳论证。如果 \(mn \neq 1\)，选择 G 的一个交换正规子群 A，它不等于 \(\{e\}\)，且其阶 a 是一个素数的幂。将归纳假设应用于 G/A 以及 \((m/a, n)\) 或 \((m, n/a)\)，视 a 整除 m 还是 n 而定。在第二种情况下，使用习题 9 从 G/A 过渡到 G。）

\begin{enumerate}
\def\labelenumi{\arabic{enumi}。}
\setcounter{enumi}{10}
\tightlist
\item
  证明 60 阶单群 \(A_{5}\) 不包含 15 阶子群。
\end{enumerate}

¶ 12. 设 G 是一个幂零群，H 是 G 中有限阶元素所成的集合。

\begin{enumerate}
\def\labelenumi{(\alph{enumi})}
\item
  证明 H 是 G 的子群。
\item
  证明 H 的每个有限子集生成一个有限子群。
\item
  证明 H 中阶互素的两个元素交换。
\end{enumerate}

\begin{enumerate}
\def\labelenumi{\arabic{enumi}。}
\setcounter{enumi}{12}
\item
  设 G 是幂零群。证明 G 是有限的（相应地，可数的）当且仅当 G/D(G) 是有限的（相应地，可数的）。
\item
  设 G 是有限群。假设对所有 \(x, y \in G\)，由 \(\{x, y\}\) 生成的 G 的子群是幂零的。证明 G 是幂零的（将假设应用于阶为 \(p_{1}^{n_{1}}\) 和 \(p_{2}^{n_{2}}\) 的 x 和 y，其中 \(p_{1}\) 和 \(p_{2}\) 是不同的素数）。
\item
  设 G 是一个群，X 是 G 的一个生成子集。
\end{enumerate}

\begin{enumerate}
\def\labelenumi{(\alph{enumi})}
\item
  记 \(X_{1} = X\)；对于 \(n \geqslant 2\)，归纳地定义 \(X_{n}\) 为换位子 \((x, y)\)（\(x \in X\)，\(y \in X_{n-1}\)）所成的集合。证明 \(X_{n} = \{e\}\) 当且仅当 G 是类 \(\leqslant n\) 的幂零群。（对 n 进行归纳论证。如果 \(X_{n} = \{e\}\)，证明由 \(X_{n-1}\) 生成的 G 的子群 H 包含在 G 的中心内，并将归纳假设应用于 G/H。）
\item
  证明 \(\mathbf{C}^n (\mathbf{G})\) 是由 \(\mathbf{G}\) 生成的 \(\mathbf{X}_n\) 的正规子群。由此推出 \(\mathbf{X}_n\) 在 \(\mathbf{C}^n (\mathbf{G}) / \mathbf{C}^{n + 1}(\mathbf{G})\) 中的像生成群 \(\mathbf{C}^n (\mathbf{G}) / \mathbf{C}^{n + 1}(\mathbf{G})\)。
\item
  取 \(\mathbf{G} = \mathfrak{S}_m\)（\(m \geqslant 4\)）且 \(\mathbf{X} = \{s, t\}\)，其中 \(s\) 是一个对换，而 \(t\) 是一个阶为 \(m\) 的轮换。证明 \(\mathbf{X}_2\) 不生成群 \(\mathbf{C}^2(\mathbf{G}) = \mathfrak{A}_m\)。
\end{enumerate}

\begin{enumerate}
\def\labelenumi{\arabic{enumi}。}
\setcounter{enumi}{15}
\tightlist
\item
  *设 V 是交换域 k 上的向量空间，且
\end{enumerate}

\[
\mathbf {W} = \bigwedge^ {2} \mathbf {V}
\]

（参见 III，§ 7，第 1 号。）在 \(\mathbf{E} = \mathbf{V} \times \mathbf{W}\) 上，由以下公式定义了一个合成律

\[
(v _ {1}, w _ {1}). (v _ {2}, w _ {2}) = (v _ {1} + v _ {2}, w _ {1} + w _ {2} + v _ {1} \wedge v _ {2}).
\]

\begin{enumerate}
\def\labelenumi{(\alph{enumi})}
\item
  证明这个合成律赋予 \(\mathbf{E}\) 一个群结构，即 \(\mathbf{V}\) 的一个以 \(\mathbf{W}\) 为核的中心扩张。
\item
  假设 \(k\) 的特征不等于 2。证明导群 \(D(E)\) 是由元素 \((0, w)\)（\(w \in W\)）组成的集合；它是同构于 \(W\) 的一个群。证明：若 \(\dim(V) \geqslant 4\)，则 \(D(E)\) 中存在不是换位子的元素。*
\end{enumerate}

\begin{enumerate}
\def\labelenumi{\arabic{enumi}。}
\setcounter{enumi}{16}
\tightlist
\item
  设 \(G\) 是一个有限交换群。证明以下条件的等价性：
\end{enumerate}

\begin{enumerate}
\def\labelenumi{(\alph{enumi})}
\item
  G 是循环的。
\item
  G 的每个西罗子群都是循环群。
\item
  对于每个素数 \(p\)，元素 \(x \in \mathbf{G}\) 中满足 \(x^p = e\) 者的个数 \(\leqslant p\)。
\end{enumerate}

\begin{enumerate}
\def\labelenumi{\arabic{enumi}。}
\setcounter{enumi}{17}
\tightlist
\item
  设 \(n\) 为整数 \(\geqslant 1\)，并设 \(G\) 是一个以加法记号书写的交换群，满足 \(nx = 0\) 对所有 \(x \in G\) 成立。设 \(H\) 是 \(G\) 的一个子群，\(f\) 是 \(H\) 到 \(\mathbf{Z} / n\mathbf{Z}\) 的一个同态。
\end{enumerate}

\begin{enumerate}
\def\labelenumi{(\alph{enumi})}
\item
  设 \(x \in \mathbf{G}\)，并设 \(q\) 是 \(x\) 在 \(\mathbf{G} / \mathbf{H}\) 中的像的阶；于是 \(qx \in \mathbf{H}\)。证明存在 \(\alpha \in \mathbf{Z} / n\mathbf{Z}\) 使得 \(q\alpha = f(x)\)。由此推出，\(f\) 可以扩张到 \(\mathbf{G}\) 的由 \(\mathbf{H}\) 和 \(x\) 生成的子群上。
\item
  证明 \(f\) 可以扩张为 \(\mathbf{G}\) 到 \(\mathbf{Z} / n\mathbf{Z}\) 的一个同态（使用佐恩引理）。
\end{enumerate}

\begin{enumerate}
\def\labelenumi{\arabic{enumi}。}
\setcounter{enumi}{18}
\tightlist
\item
  设 G 是一个有限交换群。
\end{enumerate}

\begin{enumerate}
\def\labelenumi{(\alph{enumi})}
\item
  证明存在一个元素 \(x \in G\)，其阶 n 是 G 中诸元素之阶的倍数。（利用 G 分解为 p-群的乘积。）证明由这样一个元素生成的子群 H 是 G 的一个直因子（应用上面的习题构造一个同态 \(G \to H\)，它在 H 上的限制为恒等映射）。
\item
  证明 G 是循环群的直积。（通过对 Card(G) 进行归纳并使用 (a) 来论证。）（参见第七章，第 4 节，第 6 号。）
\end{enumerate}

\begin{enumerate}
\def\labelenumi{\arabic{enumi}。}
\setcounter{enumi}{19}
\tightlist
\item
  设 \(G\) 是一个以加法记号书写的有限交换群。令 \(x = \sum_{g \in G} g\)，并设 \(G_2\) 为 \(G\) 的由元素 \(g\)（满足 \(2g = 0\)）所组成的子群。
\end{enumerate}

\begin{enumerate}
\def\labelenumi{(\alph{enumi})}
\item
  证明 \(x = \sum_{g \in G_2} g\)。
\item
  证明 \(x = 0\)（若 \(\operatorname{Card}(\mathbf{G}_2) \neq 2\)）。若 \(\operatorname{Card}(\mathbf{G}_2) = 2\)，证明 \(x\) 是 \(\mathbf{G}_2\) 中唯一的非零元素。
\item
  证明 \(\operatorname{Card}(\mathbf{G}_2) = 2\) 当且仅当 \(\mathbf{G}\) 的西罗 2-子群是循环的且 \(\neq 0\)。
\item
  *设 \(p\) 是一个素数。证明
\end{enumerate}

\[
(p - 1)! \equiv - 1 \pmod {p}.
\]

（将 (b) 应用于域 \(\mathbf{Z} / p\mathbf{Z}\) 的乘法群。）

\begin{enumerate}
\def\labelenumi{\arabic{enumi}。}
\setcounter{enumi}{20}
\tightlist
\item
  设 \(p\) 和 \(q\) 是两个素数，满足 \(p > q\)。
\end{enumerate}

\begin{enumerate}
\def\labelenumi{(\alph{enumi})}
\item
  证明每个阶为 \(pq\) 的有限群 G 都是 \(\mathbf{Z}/q\mathbf{Z}\) 的一个以 \(\mathbf{Z}/p\mathbf{Z}\) 为核的扩张（注意 G 的每个西罗 \(p\)-子群都是正规子群）。若进一步 \(p \not\equiv 1 \pmod{q}\)，证明 G 是循环群。
\item
  由此推出，一个群的中心的指数不能为 69。
\item
  证明：若 \(p \equiv 1 \pmod{q}\)，则存在一个阶为 \(pq\) 的群，其中心为 \(\{e\}\)。
\item
  证明每个 6 阶非交换群都同构于 \(\mathfrak{S}_3\)。
\end{enumerate}

\begin{enumerate}
\def\labelenumi{\arabic{enumi}。}
\setcounter{enumi}{21}
\tightlist
\item
  设 G 是一个阶为 n \textgreater{} 1 的有限群，并设 p 为整除 n 的最小素数。设 P 为 G 的一个西罗 p-子群，N 为其正规化子。证明：若 P 是循环群，则 P 包含在 N 的中心内（证明 N/P 的阶与 P 的自同构群的阶互素）。
\end{enumerate}

¶ 23. 设 σ 是群 G 的一个自同构。

\begin{enumerate}
\def\labelenumi{(\alph{enumi})}
\item
  证明以下条件的等价性：
\item
  \(\sigma(x) = x\) 蕴含 \(x = e\)。
\end{enumerate}

\begin{enumerate}
\def\labelenumi{(\roman{enumi})}
\setcounter{enumi}{1}
\tightlist
\item
  映射 \(x \mapsto x^{-1}\sigma(x)\) 是单射。
\end{enumerate}

当这些条件满足时，称 \(\sigma\)（通过语言的滥用）为无不动点的。

\begin{enumerate}
\def\labelenumi{(\alph{enumi})}
\setcounter{enumi}{1}
\item
  假设 G 是有限的，且 \(\sigma\) 无不动点。映射 \(x \mapsto x^{-1}\sigma(x)\) 因而是双射。若 H 是 G 的一个在 \(\sigma\) 下稳定的正规子群，证明由 \(\sigma\) 定义的 G/H 的自同构是无不动点的。
\item
  在 (b) 的假设下，设 p 为一个素数。证明存在 G 的一个在 \(\sigma\) 下稳定的西罗 p-子群 P。（如果 \(P_{0}\) 是一个西罗 p-子群，则存在 \(y \in G\) 使得 \(\sigma(P_{0}) = yP_{0}y^{-1}\)；把 \(y^{-1}\) 写成 \(x^{-1}\sigma(x)\) 的形式，并取 \(P = xP_{0}x^{-1}\)。）证明这样的子群是唯一的，并且包含 G 中在 \(\sigma\) 下稳定的每个 p-子群。
\item
  进一步假设 \(\sigma\) 的阶为 2。证明 \(\sigma(g) = g^{-1}\) 对所有 \(g \in G\) 成立（把 \(g\) 写成 \(x^{-1}\sigma(x)\) 的形式，其中 \(x \in G\)）。由此推出 \(G\) 是交换的且具有奇数阶。
\end{enumerate}

\begin{enumerate}
\def\labelenumi{\arabic{enumi}。}
\setcounter{enumi}{23}
\item
  设 G 为有限群，H 为 G 的西罗子群，N 为 H 的正规化子。设 \(X_{1}, X_{2}\) 是 H 的中心的两个子集，且 \(s \in G\) 满足 \(sX_{1}s^{-1} = X_{2}\)。证明存在 \(n \in N\) 使得 \(nxn^{-1} = sxs^{-1}\) 对所有 \(x \in X_{1}\) 成立（将西罗子群共轭定理应用于 \(X_{1}\) 的中心化子）。由此推出，H 的两个中心元素在 G 中共轭当且仅当它们在 N 中共轭。
\item
  设 \(\phi: G \to G'\) 是有限群之间的一个满同态；设 \(P\) 是一个西罗 \(p\)-子群（属于 \(G\)）。
\end{enumerate}

\begin{enumerate}
\def\labelenumi{(\alph{enumi})}
\item
  设 \(\mathbf{P}_1\) 是一个西罗 \(p\)-子群（\(\mathbf{G}\) 的），满足 \(\phi(\mathbf{P}) = \phi(\mathbf{P}_1)\)。证明存在一个元素 \(x\)（属于 \(\phi\) 的核），使得 \(\mathbf{P}_1 = x\mathbf{P}x^{-1}\)。
\item
  设 N（相应地 N'）为 P（相应地 \(\phi(\mathrm{P})\)）在 G（相应地 G'）中的正规化子。证明 \(\phi(\mathrm{N}) = \mathrm{N}'\)。特别地，如果 \(G'\) 的阶不被 p 整除，则 \(\phi(\mathrm{P}) = \{e\}\) 且 \(\phi(\mathrm{N}) = \mathrm{G}'\)。
\end{enumerate}

¶ 26. 设 G 是一个有限群。若存在由 G 的正规子群组成的一个合成列 \((\mathbf{G}_{i})_{0\leqslant i\leqslant n}\)，使得商群 \(G_{i}/G_{i+1}\) 都是循环群，则称 G 为超可解群。

\begin{enumerate}
\def\labelenumi{(\alph{enumi})}
\item
  证明超可解群的每个子群、每个商群以及每个有限乘积都是超可解的。
\item
  证明幂零群 \(\Rightarrow\) 超可解群 \(\Rightarrow\) 可解群，并且相反的蕴含关系不成立。
\item
  假设 G 是超可解的。证明：若 \(G \neq \{1\}\)，则存在 G 的一个正规子群 C，它是素数阶循环群。证明 G 的导群 \(D(G)\) 是幂零的。证明 G 的每个不同于 G 的极大子群都有素数指数。
\item
  一个有限群 G 称为双循环群，如果存在 G 的两个元素 a、b，使得若 \(C_{a}\) 和 \(C_{b}\) 分别表示由 a 和 b 生成的循环子群，则 \(G = C_{a}C_{b} = C_{b}C_{a}\)。证明每个双循环群都是超可解群。（将其归结为证明 G 中存在一个正规循环子群 \(N \neq \{1\}\)。可以只注意 \(C_{a} \cap C_{b} = \{1\}\) 的情形；令 \(m \geqslant n > 1\) 分别为 a 和 b 的阶；通过考虑元素 \(b^{-1}a^{k}\)（其中 \(0 \leqslant k \leqslant m - 1\)），证明存在一个整数 s 使得 \(a^{s} \neq 1\) 且 \(b^{-1}a^{s}b \in C_{a}\)；具有此性质的诸 \(a^{s}\) 所成的集合即为所求的循环子群 N。）
\end{enumerate}

\begin{enumerate}
\def\labelenumi{\arabic{enumi}。}
\setcounter{enumi}{26}
\tightlist
\item
  设 S 为一个有限群。若 S 是单群且非交换，并且 S 的每个不同于 S 的子群都是可解群，则称 S 为极小单群。
\end{enumerate}

\begin{enumerate}
\def\labelenumi{(\alph{enumi})}
\item
  证明交错群 \(\mathfrak{A}_n\) 是极小单群当且仅当 \(n = 5\)。
\item
  设 G 为有限群。证明：若 G 不可解，则存在 G 的两个子群 H 和 K，使得 H 在 K 中正规，且 K/H 为极小单群。
\end{enumerate}

¶ 28. 设 G 为有限群，p 为素数。元素 \(s \in G\) 称为 p-幂单（相应地，p-正则），如果其阶为 p 的幂（相应地，不被 p 整除）。

\begin{enumerate}
\def\labelenumi{(\alph{enumi})}
\item
  设 \(x \in G\)。证明存在唯一的有序对 \((u, t)\)（由 \(G\) 的元素组成），满足以下条件：\(u\) 是 \(p\)-幂单的，\(t\) 是 \(p\)-正则的，\(x = ut = tu\)。（首先考虑 \(G\) 是由 \(x\) 生成的循环群的情形。）
\item
  设 P 为 G 的西罗 p-子群，C 为其中心化子，E 为 G 的 p-正则元素所成的集合。证明
\end{enumerate}

\[
\operatorname{Card} (\mathbf {E}) \equiv \operatorname{Card} (\mathbf {E} \cap \mathbf {C}) \pmod {p}.
\]

由此推出 \(\operatorname{Card}(\mathbf{E}) \not\equiv 0 (\bmod p)\)。（通过对 \(\operatorname{Card}(\mathbf{G})\) 进行归纳论证，并将其归结为 \(\mathbf{C} = \mathbf{G}\) 的情形；然后利用 \((a)\) 证明 \(\operatorname{Card}(\mathbf{E}) = (\mathbf{G}:\mathbf{P}).\)）

¶ 29. 设 G 为偶数阶有限群，H 为 G 的西罗 2-子群。对所有 \(s \in G\)，设 \(\varepsilon(s)\) 为 G 的置换 \(x \mapsto sx\) 的符号。

\begin{enumerate}
\def\labelenumi{(\alph{enumi})}
\item
  设 \(s \in \mathbf{H}\)。证明 \(\varepsilon(s) = -1\) 当且仅当 \(s\) 生成 \(\mathbf{H}\)。
\item
  证明 \(\varepsilon\) 是满射当且仅当 H 是循环群。
\item
  假设 H 是循环群。证明存在且仅存在一个 G 的正规子群 D，使得 G 是 H 与 D 的半直积。（通过对 H 的阶进行归纳论证。）证明 H 在 G 中的正规化子 N 是 H 与 \(N \cap D\) 的直积。
\item
  证明一个非交换单群的阶要么是奇数†，要么能被 4 整除。
\end{enumerate}

¶ 30. 设 p 为素数，r 和 t 为整数 \(\geqslant0\)，S 是一个阶为 \(p^{r+t}\) 的循环群，T 是 S 的阶为 \(p^{t}\) 的子群。

\begin{enumerate}
\def\labelenumi{(\alph{enumi})}
\tightlist
\item
  若 S 作用于有限集合 E 上，证明
\end{enumerate}

\[
\operatorname{Card} (\mathbf {E}) \equiv \operatorname{Card} (\mathbf {E} ^ {\mathrm{T}}) \pmod {p ^ {r + 1}},
\]

其中 \(E^{T}\) 表示 E 中在 T 下不变的元素所成的集合。

\begin{enumerate}
\def\labelenumi{(\alph{enumi})}
\setcounter{enumi}{1}
\tightlist
\item
  设 \(m, n \geqslant 0\)。证明
\end{enumerate}

\[
\binom {p ^ {r + t} m} {p ^ {t} n} \equiv \binom {p ^ {r} m} {n} \pmod {p ^ {r + 1}}.
\]

（设 M 是一个具有 m 个元素的集合，E 是 S × M 的具有 \(p^{t}n\) 个元素的子集所成的集合。定义 S 在 S × M 上的一个作用，使得存在一个从 \(E^{T}\) 到 (S/T) × M 的具有 n 个元素的子集所成的集合上的双射；应用 (a)。）

\begin{enumerate}
\def\labelenumi{\arabic{enumi}。}
\setcounter{enumi}{30}
\tightlist
\item
  设 G 为有限群，S 为 G 的西罗 p-子群。设 r, t 为整数 \(\geqslant0\)，使得 \(\operatorname{Card}(S)=p^{r+t}\)。设 E 为 G 的阶为 \(p^{t}\) 的子群所成的集合。
\end{enumerate}

\begin{enumerate}
\def\labelenumi{(\alph{enumi})}
\item
  假设 S 是循环的。证明 \(\operatorname{Card}(\mathbf{E}) \equiv 1 \pmod{p^{r+1}}\)。（令 S 通过共轭作用在 E 上，并利用习题 30。）由此推出，若 S 的阶为 \(p^{t}\) 的子群在 G 中不是正规的，则 G 中至少有 \(1 + p^{r+1}\) 个西罗 p-子群。
\item
  证明 \(\operatorname{Card}(\mathbf{E}) \equiv 1 \pmod{p}\)，即使 S 不是循环的。（设 G 通过平移作用于 G 的由具有 \(p^{t}\) 个元素的子集组成的集合 F 上。证明 E 的元素给出不同的轨道，每个轨道有 \((\mathrm{G}: \mathrm{S})p^{r}\) 个元素，并且所有其他轨道的元素个数都能被 \(p^{r+1}\) 整除。应用习题 30。）
\end{enumerate}

¶ 32. 设 p 为一个素数，G 为一个 p-群。设 \(\mathrm{G}^{*} = \mathrm{G}^{p}\mathrm{D}(\mathrm{G})\) 为 G 的由 \(\mathrm{D}(\mathrm{G})\) 及诸 \(x^{p}, x \in G\) 所生成的子群。

\begin{enumerate}
\def\labelenumi{(\alph{enumi})}
\item
  证明 \(\mathbf{G}^*\) 是 \(\mathbf{G}\) 到 \(\mathbf{Z} / p\mathbf{Z}\) 的诸同态的核的交。
\item
  证明 G 的子集 S 生成 G 当且仅当 S 在 \(\mathbf{G} / \mathbf{G}^*\) 中的像生成 \(\mathbf{G} / \mathbf{G}^*\)。由此推出，若 \((\mathbf{G}:\mathbf{G}^{*}) = p^{n}\)，则整数 \(n\) 是 G 的一个生成 G 的子集所需的最少元素个数；特别地，G 是循环群当且仅当 \(n\leqslant 1\)。
\item
  设 \(u\) 是 \(G\) 的一个阶与 \(p\) 互素的自同构，并设 \(\bar{u}\) 为 \(G / G^*\) 的对应自同构。证明 \(\bar{u} = 1\) 蕴含 \(u = 1\)。
\end{enumerate}

¶ 33. 设 G 是一个 p-群，A 是它的自同构群。

\begin{enumerate}
\def\labelenumi{(\alph{enumi})}
\tightlist
\item
  设 \((\mathbf{G}_i)_{0\leqslant i\leqslant n}\) 为 \(\mathbf{G}\) 的一个合成列，使得 \((\mathbf{G}_i:\mathbf{G}_{i + 1}) = p\) 对于
\end{enumerate}

\(0 \leqslant i \leqslant n - 1\) 成立。设 \(P\) 为 \(A\) 的由满足以下条件的自同构 \(u\) 组成的子群：对所有 \(i\) 以及所有 \(x \in G_i\)，\(u(x)x^{-1} \in G_{i+1}\)。证明，若一个元素 \(u \in P\) 的阶与 \(p\) 互素，则 \(u = 1\)（对 \(n\) 进行归纳论证）。由此推出 \(P\) 是一个 \(p\)-群。

\begin{enumerate}
\def\labelenumi{(\alph{enumi})}
\setcounter{enumi}{1}
\tightlist
\item
  反之，设 P 为 A 的一个 \(p\)-子群。证明存在 G 的一个在 P 下稳定的合成列 \((\mathbf{G}_i)_{0 \leqslant i \leqslant n}\)，使得 \((\mathbf{G}_i: \mathbf{G}_{i+1}) = p\) 对于 \(0 \leqslant i \leqslant n-1\) 成立；证明 \(\mathbf{G}_i\) 可选取为在 G 中正规。
\end{enumerate}

\begin{enumerate}
\def\labelenumi{\arabic{enumi}。}
\setcounter{enumi}{33}
\item
  设 p 为一个素数。证明每个阶为 \(p^{2}\) 的群都是交换的。
\item
  设 G 为一个群，使得其导群 \(D = D(G)\) 包含在 G 的中心 C 中。
\end{enumerate}

\begin{enumerate}
\def\labelenumi{(\alph{enumi})}
\tightlist
\item
  证明映射 \((x,y)\mapsto (x,y)\) 诱导一个映射
\end{enumerate}

\[
\phi \colon (\mathbf {G} / \mathbf {C}) \times (\mathbf {G} / \mathbf {C}) \rightarrow \mathbf {D}
\]

使得，将这些群写成加法群：

\[
\begin{array}{l l} \phi (\alpha + \beta , \gamma) = \phi (\alpha , \gamma) + \phi (\beta , \gamma) & \quad \phi (\alpha , \beta) = - \phi (\beta , \alpha) \\ \phi (\alpha , \beta + \gamma) = \phi (\alpha , \beta) + \phi (\alpha , \gamma) & \quad \phi (\alpha , \alpha) = 0 \end{array}
\]

对所有 \(\alpha, \beta, \gamma \in \mathbf{G} / \mathbf{C}\) 成立。

\begin{enumerate}
\def\labelenumi{(\alph{enumi})}
\setcounter{enumi}{1}
\tightlist
\item
  证明：对每个整数 n，
\end{enumerate}

\[
(y x) ^ {n} = y ^ {n} x ^ {n} (x, y) ^ {n (n - 1) / 2}
\]

对于 \(x, y \in \mathbf{G}\) 成立。由此推出，若 \(n\) 为奇数且 \(d^n = 1\) 对所有 \(d \in \mathbf{D}\) 成立，则映射 \(x \mapsto x^n\) 诱导一个同态 \(\theta: \mathbf{G} / \mathbf{D} \to \mathbf{C}\)。

¶ 36. 设 p 为一个素数，且 G 为一个阶为 \(p^{3}\) 的非交换群。

\begin{enumerate}
\def\labelenumi{(\alph{enumi})}
\item
  证明，使用前一习题的记号，\(\mathbf{C} = \mathbf{D}\)，并且 \(\mathbf{G} / \mathbf{D}\) 同构于两个阶为 \(p\) 的群的乘积。
\item
  假设 p 是奇数，且前一习题中的同态 \(\theta\) 非零。证明 G 是一个 p 阶群与一个阶为 \(p^{2}\) 的循环群的半直积。证明存在生成 G 的元素 x, y，使得
\end{enumerate}

\[
x ^ {p ^ {2}} = 1, \quad y ^ {p} = 1, \quad (x, y) = x ^ {p},
\]

即 G 在同构意义下由该性质（以及它具有阶 \(p^{3}\) 这一事实）所刻画，并且这样的群 G 存在。

\begin{enumerate}
\def\labelenumi{(\alph{enumi})}
\setcounter{enumi}{2}
\item
  对于 \(p = 2\)，考虑与 (b) 中相同的问题，但将假设 \(\theta \neq 0\) 替换为假设 \(G\) 包含一个 2 阶的非中心元素。证明这样的群同构于二面体群 \(D_4\)（参见习题 4），也同构于 \(S_4\) 的一个西罗 2-子群。
\item
  假设 \(p\) 是奇数，且前一习题中的同态 \(\theta\) 为零。证明存在元素 \(x, y, z\)（它们生成 \(G\)），使得
\end{enumerate}

\[
(x, y) = z; \quad x ^ {p} = y ^ {p} = z ^ {p} = (x, z) = (y, z) = 1,
\]

即 G 由这一性质所刻画，并且这样的群是存在的。

*证明 G 同构于如下形式的矩阵乘法群

\[
\left( \begin{array}{c c c} 1 & a & b \\ 0 & 1 & c \\ 0 & 0 & 1 \end{array} \right)
\]

其中的系数 \(a, b, c\) 属于含 \(p\) 个元素的域。*

\begin{enumerate}
\def\labelenumi{(\alph{enumi})}
\setcounter{enumi}{4}
\tightlist
\item
  假设 \(p = 2\)，并且 \(G\) 中没有阶为 2 的非中心元素（因而它们的阶都是 4）。证明 \(G\) 由元素 \(x, y\) 生成，使得
\end{enumerate}

\[
x ^ {2} = y ^ {2} = (x, y); \quad (x, y) ^ {2} = 1,
\]

即 G 由这一性质所刻画，并且同构于四元数群（参见习题 4）。

\begin{enumerate}
\def\labelenumi{(\alph{enumi})}
\setcounter{enumi}{5}
\tightlist
\item
  *问：由形如
\end{enumerate}

\[
\left( \begin{array}{c c c} 1 & a & b \\ 0 & 1 & c \\ 0 & 0 & 1 \end{array} \right),
\]

（其中 \(a, b, c\) 遍历含 2 个元素的域）的矩阵所成的群是 \((c)\) 型还是 \((e)?\)*

\begin{enumerate}
\def\labelenumi{\arabic{enumi}。}
\setcounter{enumi}{36}
\tightlist
\item
  \begin{enumerate}
  \def\labelenumii{(\alph{enumii})}
  \tightlist
  \item
    设 G 为有限群，H 为 G 的正规子群，p 为不整除 H 的阶的素数，P 为 G 的西罗 p-子群。证明 HP 是 P 与 H 的半直积。
  \end{enumerate}
\end{enumerate}

\begin{enumerate}
\def\labelenumi{(\alph{enumi})}
\setcounter{enumi}{1}
\tightlist
\item
  我们说一个有限群 G 具有性质 (ST)，如果存在诸不同素数的一个编号 \(p_{1}, \ldots, p_{s}\)（这些素数整除阶
\end{enumerate}

\[
n = p _ {1} ^ {a _ {1}} \dots p _ {s} ^ {a _ {s}}
\]

），并且存在 G 的西罗 \(p_i\)-子群 \(P_1, \ldots, P_s\)，使得对于 \(1 \leqslant i \leqslant s\)，集合 \(G_i = P_1P_2 \ldots P_i\) 是 G 的正规子群。证明，若确实如此，则 \(G_i\) 不依赖于 \(P_i\) 的选取；它是 G 的唯一一个阶为 \(p_1^{a_1} \ldots p_i^{a_i}\) 的子群；进一步，\(G_i\) 是 \(P_i\) 与 \(G_{i-1}\) 的半直积。

\begin{enumerate}
\def\labelenumi{(\alph{enumi})}
\setcounter{enumi}{2}
\item
  证明：有限群 G 具有性质 (ST) 当且仅当存在 G 的一个西罗子群 P，使得 P 在 G 中正规，并且 G/P 具有性质 (ST)。
\item
  证明：具有性质 (ST) 的群的每个子群、每个商群以及每个中心扩张都具有性质 (ST)。
\end{enumerate}

¶ 38. 设 G 是一个阶为 \(n = p^{a}m\) 的有限群，其中 \(p\) 是一个不整除 \(m\) 的素数，并设 X 为 G 的西罗 \(p\)-子群所成的集合。设 P \ensuremath{\in} X，令 N 为 P 的正规化子，且 \(r = \text{Card}(X)\)。

\begin{enumerate}
\def\labelenumi{(\alph{enumi})}
\item
  证明 \(r = (\mathbf{G}:\mathbf{N})\)。由此推出 \(r\) 属于 \(\mathbf{R}\)（即 \(m\) 的与 1 同余（模 \(p\)）的正因子所成的集合）。特别地，若 \(\mathbf{R}\) 是 \(\{1\}\)，则 \(\mathbf{P}\) 在 \(\mathbf{G}\) 中正规。
\item
  若 \(a = 1\)，证明 \(G\) 中阶为 \(p\) 的元素的个数是 \(r(p - 1)\)。若进一步 \(N = P\)，证明 \(G\) 中阶不同于 \(p\) 的元素在个数上等于 \(m\)；若进一步 \(m\) 是一个素数 \(q\) 的幂，则推断出一个西罗 \(q\)-子群 \(Q\)（属于 \(G\)）包含 \(G\) 中所有阶不同于 \(p\) 的元素，因而在 \(G\) 中是正规的。
\item
  设 H 是一个阶为 210 的群。证明 H 包含一个阶为 105 的正规子群 G（参见习题 29）。设 p = 3（相应地，5，7）。证明：若 G 没有正规的西罗 p-子群，则群 G 至少包含 14（相应地，84，90）个阶为 p 的元素。由此得出结论：H 和 G 具有性质 (ST)（参见前一个习题）。
\item
  证明每个阶为 \(n \leqslant 100\) 的群都具有性质 (ST)，除非可能有 \(n = 24, 36, 48, 60, 72, 96\)（方法与 (c) 相同）。
\item
  证明群 \(\mathfrak{S}_4\)（其阶为 24）不具有性质 (ST)。为 \(n = 48, 60, 72, 96\) 构造类似的例子。（对于 \(n = 36\)，见下面的习题。）
\end{enumerate}

¶ 39. 设 G 为有限群，\(n = \text{Card}(G)\)，\(p\) 为素数，X 为 G 的西罗 \(p\)-子群之集，且 \(r = \text{Card}(X)\)。则 G 通过共轭传递地作用在 X 上；设 \(\phi: G \to \mathfrak{S}_{X} \approx \mathfrak{S}_r\) 为相应的同态（\(A \approx B\) 表示“A 同构于 B”这一关系），K 为其核。

\begin{enumerate}
\def\labelenumi{(\alph{enumi})}
\item
  证明 K 是诸 \(P \in X\) 的正规化子的交，且 K 的西罗 p-子群是诸 \(P \in X\) 的交。由此推出，若 r \textgreater{} 1，则 K 的阶 k 整除 n/rp。
\item
  证明：若 G 没有指数为 2 的子群，则 \(\phi(G) \subset \mathfrak{A}_{X}\)；并且若 \(n = kr!\)（相应地，\(\frac{1}{2}kr!\)），则 \(\phi(G) = \mathfrak{S}_{X}\)（相应地，\(\mathfrak{A}_{X}\)）。
\item
  利用 (a)、(b) 及前一习题，证明以下事实：
\end{enumerate}

若 \(n = 12\) 并且 \(G\) 的西罗 3-子群不正规，则 \(G \approx \mathfrak{A}_4\)。

若 \(n = 24\) 且 \(\mathbf{G}\) 不具有性质 (ST)，则 \(\mathrm{G} \approx \mathfrak{S}_4\)。

若 n = 36，则 G 具有性质 (ST)。（证明：若 G 的一个西罗 3-子群不是正规的，则 G 包含一个 3 阶子群 K，使得 \(G/K \approx A_{4}\) 且 K 是中心的。）

若 \(n = 48\) 且 \(G\) 不具有性质 (ST)，则 \(G\) 包含一个阶为 2 的正规子群 \(K\)，使得 \(G / K \approx S_4\)。

若 \(n = 60\) 且 \(G\) 不具有性质 (ST)，则 \(G \approx \mathfrak{A}_5\)。（证明这样的群 \(G\) 没有 5 阶正规子群，没有 2 阶正规子群，也没有指数为 2 的正规子群。由此推出 \(G\) 同构于 \(\mathfrak{S}_6\) 的一个子群，并且，按照 § 5 习题 25 中的论证方式，\(G \approx \mathfrak{A}_5\)。）

若 \(n = 72\) 并且 \(G\) 的西罗 3-子群不正规，则 \(G\) 包含一个正规子群 \(K\)，使得 \(G / K\) 同构于 \(\mathfrak{S}_4\) 或同构于 \(\mathfrak{A}_4\)。

若 n = 96 且 G 的一个西罗 2-子群不是正规的，则 G 包含一个正规子群 K，使得 G/K ≈ \(S_{3}\)。

\begin{enumerate}
\def\labelenumi{(\alph{enumi})}
\setcounter{enumi}{3}
\tightlist
\item
  证明一个阶 \(\leqslant 100\) 的群 G 是可解的，除非 \(\mathbf{G} \approx \mathfrak{A}_5\)。
\end{enumerate}

\begin{enumerate}
\def\labelenumi{\arabic{enumi}。}
\setcounter{enumi}{39}
\tightlist
\item
  设 \(G_{1}, G_{2}\) 是两个群，G 是 \(G_{1} \times G_{2}\) 的一个子群，满足 \(pr_{1} G = G_{1}, pr_{2} G = G_{2}\)；把 \(G_{1}\) 和 \(G_{2}\) 分别与 \(G_{1} \times \{e_{2}\}\) 和 \(\{e_{1}\} \times G_{2}\) 等同起来，则 \(H_{1} = G \cap G_{1}\) 和 \(H_{2} = G \cap G_{2}\) 是 G 的正规子群，使得 \(G/H_{1}\) 同构于 \(G_{2}\)，且 \(G/H_{2}\) 同构于 \(G_{1}\)。
\end{enumerate}

\begin{enumerate}
\def\labelenumi{(\alph{enumi})}
\item
  证明：要使 G 在 \(G_{1} \times G_{2}\) 中正规，必要且充分的是 G 包含导群 \(\mathrm{D}(\mathrm{G}_{1} \times \mathrm{G}_{2}) = \mathrm{D}(\mathrm{G}_{1}) \times \mathrm{D}(\mathrm{G}_{2})\)（证明 G 包含 \(\mathrm{D}(\mathrm{G}_{1})\)：写出 \(xzx^{-1}z^{-1} \in G\)，其中 \(x \in G_{1}\)、\(z \in G\)）。由此推出 \(G_{1}/H_{1}\) 是交换的。
\item
  若 G 在 \(G_{1} \times G_{2}\) 中正规，证明有一个 \(G_{1} \times G_{2}\) 到 \(G_{1}/H_{1}\) 上的同态被如下定义：把每个有序对 \((s_{1}, s_{2})\) 对应到 \(s_{1}t_{1}^{-1}\) 模 \(H_{1}\) 的陪集，其中 \(t_{1} \in G_{1}\) 满足 \((t_{1}, s_{2}) \in G\)。由此推出 \((\mathrm{G}_{1} \times \mathrm{G}_{2})/\mathrm{G}\) 同构于 \(\mathrm{G}/(\mathrm{H}_{1} \times \mathrm{H}_{2})\)。
\end{enumerate}

\begin{enumerate}
\def\labelenumi{\arabic{enumi}。}
\setcounter{enumi}{40}
\tightlist
\item
  设 \(x, y\) 是群 \(G\) 的两个元素。
\end{enumerate}

\begin{enumerate}
\def\labelenumi{(\alph{enumi})}
\item
  为使存在 \(a, b\)（属于 \(G\)）满足 \(bay = xab\)，必要且充分的是 \(xy^{-1}\) 是一个换位子。
\item
  为使在 G 中存在 \(2n + 1\) 个元素 \(a_1, a_2, \ldots, a_{2n+1}\)，使得
\end{enumerate}

\[
x = a _ {1} a _ {2} \dots a _ {2 n + 1} \quad \text { and } \quad y = a _ {2 n + 1} a _ {2 n} \dots a _ {1},
\]

其充要条件是 \(xy^{-1}\) 是 n 个换位子的乘积（对 n 进行归纳论证）。
% semantic-fragments: legacy-input-seam
\relax{}
